%% Research Note of the AAS (Be-star TESS variability vs WISE disk excess), AASTeX v7.0.1.
%% v3 2026-11-24: submitted as AAS84043. v4 2026-11-25: wording shortened for AAS's word count (references
%% included); no change to numbers or results. Not uploaded.
%% Provenance of every number: rnaas/numbers.md. Code: https://github.com/coracle-us/duo-transits.
\documentclass{aastex701}

\begin{document}

\title{Be-star Candidates Are More Variable in TESS When Their WISE Disk Excess Is Larger}

\author[orcid=0009-0000-7524-4470]{Jon McBee}
\affiliation{Independent researcher, USA}
\email{jon.mcbee@gmail.com}

\begin{abstract}
Be stars build and lose their decretion disks over months to years. I compare TESS light curves with the NEOWISE W1 light curves of
1,084 Be-star candidates from the catalogue of \citet{Jia2026}, sector by sector and within each
star. When the W1 disk excess is larger, the star is more variable in the optical. The
variability amplitude below 1\,c\,d$^{-1}$ rises by about 8\% per 0.1\,mag of excess and passes a
null test that swaps disk histories between stars ($p<10^{-4}$). It also rises in the 256 stars
with H$\alpha$ emission in their Gaia spectra. At higher frequencies the rise is about 3\%, and
the dominant periodic signal shows no separate response. Stars whose excess changed most carry
the trend: of 177 that changed by more than 0.4\,mag, 71\% show it. This is a population
tendency, not a rule for each star.
\end{abstract}

\keywords{\uat{Be stars}{142} --- \uat{Circumstellar disks}{235} --- \uat{Stellar pulsations}{1625} --- \uat{Time domain astronomy}{2109}}

\section{Data and Method}

\citet{Jia2026} selected 3,300 Gaia OBA candidates whose W1 brightness and W1$-$W2 colour changed
together, the signature of a changing disk; most lack spectroscopic confirmation. Of the
3,237 not flagged as eclipsing binaries, 2,541 have $\geq2$ TESS sectors within the NEOWISE
baseline \citep{Wright2010,Mainzer2014}. Each star uses one product: QLP \citep{Huang2020},
TESS-SPOC \citep{Caldwell2020}, SPOC 120-s or GSFC-ELEANOR-LITE. Each sector was binned to 30\,min and spline-detrended (5-d knots). From its amplitude
spectrum (0.25--24\,c\,d$^{-1}$) I measured the 0.25--0.8\,c\,d$^{-1}$ rms ($B_1$); the median
residual amplitude at 0.5--3\,c\,d$^{-1}$ after prewhitening (red noise); the rms of 1-d-smoothed
residuals (flicker); the strongest peak's amplitude $A_1$; and the median 15--24\,c\,d$^{-1}$
amplitude (white noise, a control).

For W1 I added later IRSA NEOWISE visits to the catalogue's and fitted a smoothing spline like \citet{Jia2026}. A sector's disk excess is the spline at mid-sector,
measured from the star's faintest W1 level. TESS and WISE view a star about three months apart, so
the interpolated excess biases $\beta$ towards zero. A sector is used only if the flanking
visits are $\leq250$\,d apart and the spline stays within 0.1\,mag of them. The main
sample keeps sectors with CROWDSAP $\geq0.9$ (SPOC products) or a contamination ratio $\leq0.1$
(others), and stars with $\geq2$ such sectors: 1,084 stars, 4,108 sectors.

For each metric $m$ I fitted $\ln m$ against disk excess with a per-star mean removed, so
$\beta$ (per mag) uses only within-star changes. Uncertainties are 68\% bootstrap intervals
over stars. Two nulls of 10,000 draws give the $p$-values: permuting sectors within each star, and a swap
giving each star another star's W1 spline at its own sector times.
Keeping trends shared by same-epoch sectors, the swap tests whether a star's own disk
history matters. Injection tests recover $\beta$ and give
calibrated nulls.

\section{Results}

The low-frequency signal rises with disk excess (Figure~\ref{fig:beta}): $\beta=0.80\,[0.70,0.91]$
for $B_1$, 0.74 for red noise and $0.84\,[0.73,0.96]$ for flicker, with $p<10^{-4}$ under both
nulls. The slopes hold without
the crowding cut ($B_1$ 0.59) or the 100 highest-leverage stars ($B_1$ 1.03), and
with straight-line W1 interpolation or TESS-sector fixed effects (every $\beta$ within 0.03). In the 256
stars Gaia DR3 classifies as Be from H$\alpha$ emission
\citep{Creevey2023}, $B_1$ gives $0.58\,[0.29,0.87]$ ($p_{\rm swap}=5\times10^{-4}$).

In post hoc narrow frequency bins, $\beta=0.71$--0.92 below 1\,c\,d$^{-1}$ and
0.31--0.37 above 5\,c\,d$^{-1}$. Relative to each star's dominant frequency $f_1$, $\beta$ is
0.74--1.17 below $0.9f_1$, $0.43\,[0.34,0.53]$ in the bin centred on $f_1$, and 0.28--0.49 above.
$A_1$ gives 0.40, but the strongest peak is coherent in some stars, stochastic in others. In the 510 stars whose $f_1$ stays within 10\% across sectors, $A_1$
gives $0.23\,[0.10,0.36]$, and 0.08 ($p\geq0.29$) where $f_1>0.8$\,c\,d$^{-1}$. The dominant pulsation shows
no separate response.

The white-noise control also rises ($\beta=0.33$), but not where photon noise dominates: in
the noisiest third of stars it is flat (0.00, $p\geq0.87$) while $B_1$ still rises (0.63,
$p<10^{-4}$). It tracks $B_1$ (within-star correlation 0.47; slope
$-0.03\pm0.09$ with $\ln B_1$ in the fit), so the band likely carries the stochastic signal's
tail and normalizing by it over-corrects. Even so, normalized $B_1$, red noise and flicker rise
(0.41--0.51, $p<10^{-3}$), though not significantly in SPOC stars alone ($p_{\rm swap}\geq0.13$).

Added disk light dilutes SPOC fluxes, which are relative to the measured brightness, but not
QLP's, which add the TIC-expected flux to a difference-image flux \citep{Huang2020}. From SAP photometry of 125 SPOC and TESS-SPOC stars, $T$
brightens by $k=0.21$ (95\% range 0.11--0.30) per mag of W1 excess, which lowers SPOC
$\beta$ by $0.921k$ if the variable light is photospheric. Corrected, the 329 SPOC stars give 0.37
($A_1$), 0.88 ($B_1$) and 0.94 (flicker) ($p_{\rm swap}\leq0.005$), close to the 755 QLP
stars (0.52, 0.86, 0.90). Removing four likely shell stars (e.g.\ TIC~55295028) changes $\beta$ by $<0.01$.

The effect is small against the scatter: within-star rank correlations are 0.09--0.20, and at
fixed excess $\ln A_1$ varies by 0.54 rms between sectors. Of the stars,
48\% change by $<0.1$\,mag between sectors and cannot show a trend; of the 177 that change
by $>0.4$\,mag, 71\% have a positive $B_1$ slope. The $p$-values are uncorrected for
multiple metrics and samples; the main results sit at the $10^{-4}$ floor.

\begin{figure}[ht!]
\plotone{rnaas_fig1.pdf}
\caption{Within-star change in (a) $\ln B_1$ (0.25--0.8\,c\,d$^{-1}$ rms) and (b) $\ln$ flicker rms
against W1 disk excess, for 4,108 sectors of 1,084 Be-star candidates, with star means
removed. Symbols: disk phase; squares: binned means; line and band: pooled
$\beta$ and 68\% interval. Plotting code was written with Claude
(Anthropic).\label{fig:beta}}
\end{figure}

\section{Discussion}

Earlier work found the pulsation spectrum changing through outbursts of single stars and small
samples \citep{Huat2009,Semaan2018}, with extra power on the low-frequency side of the frequency
groups and in stochastic variability \citep{LabadieBartz2022,LabadieBartz2025}. $B_1$ and flicker
also include days-long outburst brightenings, so a star with a growing or dense disk
plausibly has more ejection episodes per sector. I extend this to years-long disk states in
1,084 stars. \citet{Niwano2024} found the main pulsation amplitude falling as the disk grew in four
of five Be X-ray binaries; one, RX~J0440.9+4431, is in my sample, also with a negative slope.
Their statistic is not my $A_1$, but with 48\% of my stars having a negative $A_1$ slope, four
of five is unremarkable. Any extra dependence on disk phase
depends on the W1 interpolation, so I claim none. Code and per-sector tables: \url{https://github.com/coracle-us/duo-transits}.

\begin{acknowledgments}
This work uses MAST TESS light curves: SPOC
\dataset[doi:10.17909/t9-nmc8-f686]{https://doi.org/10.17909/t9-nmc8-f686}, TESS-SPOC
\dataset[doi:10.17909/t9-wpz1-8s54]{https://doi.org/10.17909/t9-wpz1-8s54}, QLP
\dataset[doi:10.17909/t9-r086-e880]{https://doi.org/10.17909/t9-r086-e880} and GSFC-ELEANOR-LITE
\dataset[doi:10.17909/j2yt-t417]{https://doi.org/10.17909/j2yt-t417}; NEOWISE-R
photometry from IRSA \dataset[doi:10.26131/IRSA144]{https://doi.org/10.26131/IRSA144};
Gaia DR3; and the \citet{Jia2026} WISE light curves (China-VO).
\end{acknowledgments}

\facilities{TESS, WISE, NEOWISE, Gaia}
\software{lightkurve \citep{Lightkurve2018}, astroquery \citep{Ginsburg2019}, astropy, scipy, matplotlib}

\section*{AI disclosure statement}
I used Anthropic's Claude, an AI assistant, under my supervision to write and run the analysis code
(including tests and a review by a separate AI agent) and to draft
this note. I checked every number against the saved outputs and every reference against arXiv or
Crossref, and take full responsibility for the content.

\begin{thebibliography}{}
\bibitem[Caldwell et al.(2020)]{Caldwell2020} Caldwell, D.~A., Tenenbaum, P., Twicken, J.~D., et al. 2020, RNAAS, 4, 201
\bibitem[Creevey et al.(2023)]{Creevey2023} Creevey, O.~L., Sordo, R., Pailler, F., et al. 2023, A\&A, 674, A26
\bibitem[Ginsburg et al.(2019)]{Ginsburg2019} Ginsburg, A., Sip{\H o}cz, B.~M., Brasseur, C.~E., et al. 2019, AJ, 157, 98
\bibitem[Huang et al.(2020)]{Huang2020} Huang, C.~X., Vanderburg, A., P\'al, A., et al. 2020, RNAAS, 4, 204
\bibitem[Huat et al.(2009)]{Huat2009} Huat, A.-L., Hubert, A.-M., Baudin, F., et al. 2009, A\&A, 506, 95
\bibitem[Jia et al.(2026)]{Jia2026} Jia, Q., Jian, M., Li, G., Cui, Q., \& Li, C. 2026, arXiv:2611.02450
\bibitem[Labadie-Bartz et al.(2022)]{LabadieBartz2022} Labadie-Bartz, J., Carciofi, A.~C., Henrique de Amorim, T., et al. 2022, AJ, 163, 226
\bibitem[Labadie-Bartz et al.(2025)]{LabadieBartz2025} Labadie-Bartz, J., Carciofi, A.~C., Rubio, A.~C., et al. 2025, A\&A, 699, A82
\bibitem[Lightkurve Collaboration et al.(2018)]{Lightkurve2018} Lightkurve Collaboration, Cardoso, J.~V.~d.~M., Hedges, C., et al. 2018, ascl:1812.013
\bibitem[Mainzer et al.(2014)]{Mainzer2014} Mainzer, A., Bauer, J., Cutri, R.~M., et al. 2014, ApJ, 792, 30
\bibitem[Niwano et al.(2024)]{Niwano2024} Niwano, M., Fausnaugh, M.~M., Lau, R.~M., et al. 2024, MNRAS, 534, 2168
\bibitem[Semaan et al.(2018)]{Semaan2018} Semaan, T., Hubert, A.-M., Zorec, J., et al. 2018, A\&A, 613, A70
\bibitem[Wright et al.(2010)]{Wright2010} Wright, E.~L., Eisenhardt, P.~R.~M., Mainzer, A.~K., et al. 2010, AJ, 140, 1868
\end{thebibliography}

\end{document}
